\documentclass[11pt,a4paper]{article}

\usepackage[utf8]{inputenc}
\usepackage[T1]{fontenc}
\usepackage[margin=2.5cm]{geometry}
\usepackage{amsmath,amssymb,amsthm}
\usepackage{booktabs}
\usepackage{array}
\usepackage{hyperref}

% ---------------------------------------------------------------------------
% Notation macros
% ---------------------------------------------------------------------------
\newcommand{\Jset}{\mathcal{J}}
\newcommand{\Oset}{\mathcal{O}}
\newcommand{\Mset}{\mathcal{M}}
\newcommand{\Fset}{\mathcal{F}}
\newcommand{\Bset}{\mathcal{B}}
\newcommand{\Eset}{\mathcal{E}}
\newcommand{\Vset}{\mathcal{V}}
\newcommand{\job}{\operatorname{job}}
\newcommand{\cat}{\,\Vert\,}
\newcommand{\R}{\mathbb{R}}

\title{Representing the Batching Constraint in a Heterogeneous\\
Graph for the Flexible Job Shop Problem}
\author{Mariana Correia}
\date{\today}

\begin{document}
\maketitle

\section{Base problem and graph}
\label{sec:base}

\subsection{Flexible Job Shop Problem}
Let $\Jset=\{1,\dots,n\}$ be the set of jobs and $\Mset=\{1,\dots,m\}$ the set of
machines. Each job $j\in\Jset$ is an ordered sequence of $n_j$ operations
$\Oset_j=(O_{j,1},\dots,O_{j,n_j})$, and $\Oset=\bigcup_{j}\Oset_j$.
For an operation $o=O_{j,k}$ define
\begin{equation}
  \job(o)=j, \qquad \kappa(o)=k,
\end{equation}
where $\kappa(o)$ is the \emph{order index} of $o$ inside its job.
Operation $o$ can be processed on any machine of its eligible set
$\Mset_o\subseteq\Mset$, with processing time $p_{o,\mu}>0$ for $\mu\in\Mset_o$
(and $p_{o,\mu}=0 \Leftrightarrow \mu\notin\Mset_o$).

\subsection{Heterogeneous graph}
Following Echeverria (2025), the state is a heterogeneous graph
$G=(\Vset,\Eset)$ with
$\Vset=\Vset_J\cup\Vset_O\cup\Vset_M$ (job, operation and machine nodes) and the
relations
\begin{itemize}
  \item $(\text{operation},\textsf{belongs},\text{job})$: $o\to\job(o)$;
  \item $(\text{operation},\textsf{prec},\text{operation})$: $O_{j,k+1}\to O_{j,k}$, plus self-loops;
  \item $(\text{operation},\textsf{exec},\text{machine})$ and its reverse: $o\leftrightarrow\mu$ for $\mu\in\Mset_o$, each carrying an attribute vector $\mathbf{a}_{o\mu}\in\R^5$;
  \item $(\text{job},\textsf{listens},\text{job})$ and $(\text{machine},\textsf{listens},\text{machine})$: fully connected.
\end{itemize}
The node features are $\mathbf{x}_j\in\R^4$, $\mathbf{x}_o\in\R^2$ and $\mathbf{x}_\mu\in\R^3$.

% ===========================================================================
\section{FJSP with batching}
\label{sec:batching}

\subsection{Families}
Let $\Fset=\{1,\dots,|\Fset|\}$ be a set of \emph{families} and
$\varphi:\Oset\to\Fset$ the \emph{family map}. Operations of the same family have
similar characteristics and may be processed together in a \emph{batch}. Let
\begin{equation}
  \Oset^f=\{o\in\Oset:\varphi(o)=f\}.
\end{equation}

\begin{table}[h]
\centering
\caption{Family parameters.}
\label{tab:family}
\begin{tabular}{>{$}l<{$} l}
\toprule
\text{Symbol} & Meaning \\
\midrule
B_f\in\mathbb{N}           & maximum number of operations in a batch of family $f$ \\
\Delta\in\mathbb{N}        & maximum order difference inside a batch ($\Delta=2$) \\
\Mset_f=\bigcap_{o\in\Oset^f}\Mset_o & machines eligible for family $f$ (may contain several machines) \\
\bar p_{f,\mu}=\max_{o\in\Oset^f}p_{o,\mu} & upper bound on the batch processing time of $f$ on $\mu$ \\
\bottomrule
\end{tabular}
\end{table}

Two rules define which operations may be processed together:
\begin{align}
  |\Oset^f\cap\Oset_j| &\le 1 && \forall f\in\Fset,\ \forall j\in\Jset,
  \tag{F1}\label{eq:F1}\\
  |\kappa(o)-\kappa(o')| &\le \Delta && \forall o,o' \text{ in the same batch},
  \tag{F2}\label{eq:F2}
\end{align}
with $\Delta=2$. Rule~\eqref{eq:F1} is a property of the instance: a family
contains at most one operation from each job. Rule~\eqref{eq:F2} is the
\emph{order difference} constraint and formalises ``similar characteristics''.
Two operations can share a batch only if their positions in their jobs' routings
differ by at most $\Delta$. A family may contain operations that are further apart
than $\Delta$, but those operations can never be in the same batch.
Rule~\eqref{eq:F2} is enforced in the model by~\eqref{eq:B2a}. The family defines
which operations are candidates for batching and the order difference restricts
which candidates are compatible. For this reason the family feature and the
batching constraint are introduced together.

\subsection{Mixed-integer formulation}
For each family $f$, let $\Bset_f=\{1,\dots,|\Oset^f|\}$ be the index set of
\emph{potential batches}, and $\Bset=\bigcup_f\Bset_f$. Every batch $b\in\Bset_f$
has family $f(b)=f$. Let $L$ be a sufficiently large constant, for example
$L=\sum_{o}\max_\mu p_{o,\mu}$.
Batching is \emph{parallel}: the operations of a batch are processed
simultaneously on one machine, and the batch lasts as long as its longest
operation. No setup time is considered.

\paragraph{Decision variables.}
\begin{itemize}
  \item $x_{o,b}\in\{0,1\}$: operation $o$ is assigned to batch $b$, defined only for $b\in\Bset_{\varphi(o)}$;
  \item $u_b\in\{0,1\}$: batch $b$ is used;
  \item $y_{b,\mu}\in\{0,1\}$: batch $b$ is processed on machine $\mu\in\Mset_{f(b)}$;
  \item $z_{b,b'}\in\{0,1\}$: batch $b$ precedes batch $b'$ if both are on the same machine;
  \item $S_b, C_b, p_b\ge 0$: start, completion and processing time of batch $b$;
  \item $S_o, C_o\ge 0$: start and completion time of operation $o$;
  \item $C_{\max}\ge 0$: makespan.
\end{itemize}

\paragraph{Model.}
\begin{align}
  \min\ & C_{\max} \label{eq:obj}\\[2pt]
  \text{s.t.}\
  & \sum_{b\in\Bset_{\varphi(o)}} x_{o,b} = 1
    && \forall o\in\Oset \tag{B1}\label{eq:B1}\\
  & x_{o,b}=0 \quad \text{if } \varphi(o)\neq f(b)
    && \forall o\in\Oset,\ b\in\Bset \tag{B2}\label{eq:B2}\\
  & x_{o,b}+x_{o',b} \le 1
    && \forall b,\ o,o'\in\Oset^{f(b)}:|\kappa(o)-\kappa(o')|>\Delta \tag{B2a}\label{eq:B2a}\\
  & \sum_{o\in\Oset^{f(b)}} x_{o,b} \le B_{f(b)}\,u_b,\qquad x_{o,b}\le u_b
    && \forall b\in\Bset,\ o\in\Oset^{f(b)} \tag{B3}\label{eq:B3}\\
  & \sum_{\mu\in\Mset_{f(b)}} y_{b,\mu} = u_b,\qquad y_{b,\mu}=0\ \ \forall\mu\notin\Mset_{f(b)}
    && \forall b\in\Bset \tag{B4}\label{eq:B4}\\
  & p_b \ge p_{o,\mu}\,\bigl(x_{o,b}+y_{b,\mu}-1\bigr)
    && \forall b,\ o\in\Oset^{f(b)},\ \mu\in\Mset_{f(b)} \tag{B4a}\label{eq:B4a}\\
  & C_b = S_b + p_b
    && \forall b\in\Bset \tag{B5}\label{eq:B5}\\
  & S_b - L(1-x_{o,b}) \le S_o \le S_b + L(1-x_{o,b})
    && \forall o,\ b\in\Bset_{\varphi(o)} \tag{B5a}\label{eq:B5a}\\
  & C_b - L(1-x_{o,b}) \le C_o \le C_b + L(1-x_{o,b})
    && \forall o,\ b\in\Bset_{\varphi(o)} \tag{B5b}\label{eq:B5b}\\
  & S_{O_{j,k+1}} \ge C_{O_{j,k}}
    && \forall j\in\Jset,\ k<n_j \tag{B6}\label{eq:B6}\\
  & S_{b'} \ge C_b - L\bigl(3-y_{b,\mu}-y_{b',\mu}-z_{b,b'}\bigr)
    && \forall b<b',\ \mu\in\Mset \tag{B7}\label{eq:B7}\\
  & S_{b} \ge C_{b'} - L\bigl(2-y_{b,\mu}-y_{b',\mu}+z_{b,b'}\bigr)
    && \forall b<b',\ \mu\in\Mset \tag{B7'}\label{eq:B7b}\\
  & C_{\max}\ge C_b
    && \forall b\in\Bset. \label{eq:cmax}
\end{align}

\paragraph{Interpretation.}
\begin{itemize}
  \item \eqref{eq:B1}--\eqref{eq:B2}: every operation belongs to exactly one batch, and every batch contains only one family.
  \item \eqref{eq:B2a}: order difference. Two operations whose order indices differ by more than $\Delta=2$ cannot be in the same batch. Every batch $b$ therefore satisfies $\max_{o\in b}\kappa(o)-\min_{o\in b}\kappa(o)\le\Delta$.
  \item \eqref{eq:B3}: batch capacity. Unused batches are empty.
  \item \eqref{eq:B4}: a family may have several eligible machines ($|\Mset_f|\ge1$), and every used batch is assigned to exactly one of them. This is the flexibility of the FJSP, lifted to the batch level.
  \item \eqref{eq:B4a}: parallel batching. The batch processing time is the longest processing time among its members on the selected machine, $p_b=\max_{o\in b}p_{o,\mu}$. Minimising $C_{\max}$ drives $p_b$ down to this maximum.
  \item \eqref{eq:B5}--\eqref{eq:B5b}: all operations of a batch start together and finish together when the longest one finishes.
  \item \eqref{eq:B6}: routing precedence inside each job.
  \item \eqref{eq:B7}--\eqref{eq:B7b}: batches that share a machine do not overlap.
\end{itemize}
By \eqref{eq:F1} and \eqref{eq:B2}, a batch never contains two operations of the same
job, so the simultaneity constraints \eqref{eq:B5a}--\eqref{eq:B5b} never conflict
with the precedence constraint \eqref{eq:B6} inside a batch. Without batching
($B_f=1$ for all $f$), the model reduces to the classical FJSP.

% ===========================================================================
\section{Family feature on the existing nodes}
\label{sec:familyfeature}

The family index is a label that only has meaning inside one instance: family~3 of
one instance has nothing in common with family~3 of another. Feeding the index to
the network, as a number or as a learned embedding, would therefore carry no
transferable information. Which operations share a family is instead encoded
\emph{structurally}, by the family node (Representation~A) or by the batch edges
(Representation~B). The existing node features are extended with numeric
descriptors of the family:
\begin{align}
  \text{Operation node:}\quad &
  \mathbf{x}'_o = \Bigl[\ \mathbf{x}_o \cat \mathbb{1}[\varphi(o)\neq\varnothing] \cat \tfrac{\kappa(o)}{n_{\job(o)}}\ \Bigr],
  \label{eq:xop}\\
  \text{Job node:}\quad &
  \mathbf{x}'_j(t) = \Bigl[\ \mathbf{x}_j(t) \cat \mathbb{1}[\varphi(o_j^t)\neq\varnothing] \cat B_{\varphi(o_j^t)} \cat \pi_j(t)\ \Bigr],
  \label{eq:xjob}
\end{align}
where $\mathbb{1}[\varphi(o)\neq\varnothing]$ indicates that $o$ belongs to a
family, $o_j^t$ is the next unscheduled operation of job $j$ at decision step $t$,
$B_{\varphi(o_j^t)}$ is its capacity ($1$ if it has no family), and
\begin{equation}
  \pi_j(t)=\min\Bigl(B_{\varphi(o_j^t)}-1,\ \bigl|\{j'\neq j:\ \varphi(o_{j'}^t)=\varphi(o_j^t),\
  |\kappa(o_{j'}^t)-\kappa(o_j^t)|\le\Delta,\ \Mset_{o_{j'}^t}\cap\Mset_{o_j^t}\neq\emptyset,\
  r_{j'}\le s_j\}\bigr|\Bigr)
\end{equation}
is the number of other jobs that would join the batch if job $j$ were dispatched at
its earliest start $s_j$ ($r_{j'}$ is the ready time of job $j'$). It follows exactly
the batch completion rule of Section~\ref{sec:mdp}. The operation's family is static. The job's family is dynamic and changes
as the job advances through its routing. The normalised order index in
\eqref{eq:xop} gives the network the information it needs to reason about the
order difference~\eqref{eq:F2}.

% ===========================================================================
\section{Representation A: family (batch) node type}
\label{sec:repA}

A fourth node type is added:
\begin{equation}
  \Vset=\Vset_J\cup\Vset_O\cup\Vset_M\cup\Vset_F,\qquad |\Vset_F|=|\Fset|.
\end{equation}
Node $f\in\Vset_F$ represents the open batch of family $f$.

\paragraph{Node features.}
\begin{equation}
  \mathbf{x}_f(t)=\Bigl[\ B_f \cat \bar p_f \cat \tfrac{|\Mset_f|}{m} \cat \tfrac{n_f(t)}{B_f} \cat \tfrac{\kappa^{\max}_f(t)-\kappa^{\min}_f(t)}{\Delta} \cat \tfrac{r_f(t)}{|\Oset^f|}\ \Bigr],
  \label{eq:xfam}
\end{equation}
with
\begin{equation}
  \bar p_f=\frac{1}{|\Mset_f|}\sum_{\mu\in\Mset_f}\bar p_{f,\mu},\qquad
  R_f(t)=\{o_j^t\in\Oset^f:\ r_j\le t_f\},\qquad
  n_f(t)=|R_f(t)|,
\end{equation}
where $t_f=\min_{j:\,o_j^t\in\Oset^f} s_j$ is the earliest time a batch of $f$
can start, so $R_f(t)$ holds the members that could share that batch under the
completion rule of Section~\ref{sec:mdp},
where $\kappa^{\min}_f(t),\kappa^{\max}_f(t)$ are the smallest and largest order
indices in $R_f(t)$ and $r_f(t)$ is the number of members not yet scheduled. The
components are the maximum number of operations, the batch processing time (the
longest processing time, averaged over the eligible machines), the eligible
machines, the current fill level, the current order spread and the fraction of
the family still to be scheduled. The family is identified by the node itself,
so no identifier feature is needed (Section~\ref{sec:familyfeature}). The fill
level $n_f(t)/B_f$ is dynamic, and $1-n_f(t)/B_f$ is the remaining capacity. An
order spread above $1$ means that not all ready operations can share one batch
because of~\eqref{eq:F2}.

\paragraph{New relations.}
\begin{align}
  \Eset_{OF} &= \{(o,\varphi(o)) : o\in\Oset\}
  && (\text{operation},\textsf{in\_family},\text{family}) \text{ and reverse } \textsf{has},\\
  \Eset_{FM} &= \{(f,\mu) : f\in\Fset,\ \mu\in\Mset_f\}
  && (\text{family},\textsf{batch\_exec},\text{machine}) \text{ and reverse},
\end{align}
plus a self-loop $(\text{family},\textsf{self},\text{family})$ on every family node,
so that a family node keeps its own state after all its operations have been
scheduled. When the last member of a family is scheduled, its
\textsf{batch\_exec} edges are removed.
Each $(f,\mu)\in\Eset_{FM}$ carries the machine-specific batch information. It uses
the same dimension as the existing \textsf{exec} edges:
\begin{equation}
  \mathbf{a}_{f\mu}=\Bigl[\ \bar p_{f,\mu},\ \frac{\bar p_{f,\mu}}{\sum_{\mu'\in\Mset_f}\bar p_{f,\mu'}},\ 0,\ 0,\ 0\ \Bigr]\in\R^5,
\end{equation}

\paragraph{Size.} $|\Eset_{OF}|=|\Oset|$ and $|\Eset_{FM}|=\sum_f|\Mset_f|\le|\Fset|\,m$. The graph grows linearly.

% ===========================================================================
\section{Representation B: batch edge type}
\label{sec:repB}

No node type is added. Batch compatibility is encoded directly between operations
by a new relation $(\text{operation},\textsf{batch},\text{operation})$:
\begin{equation}
  \Eset_{OO}^{\mathrm{batch}}=\Bigl\{(o,o')\in\Oset^2:\
  \varphi(o)=\varphi(o'),\ \job(o)\neq\job(o'),\
  |\kappa(o)-\kappa(o')|\le\Delta,\
  \Mset_o\cap\Mset_{o'}\neq\emptyset\Bigr\}.
  \label{eq:Ebatch}
\end{equation}
The relation is symmetric. It has no self-loops, because every operation already
has one in the \textsf{prec} relation. The condition $|\kappa(o)-\kappa(o')|\le\Delta$ applies the order
difference~\eqref{eq:F2} directly in the graph. Operations of the same family that
are more than $\Delta=2$ positions apart are not connected, so the network never
receives a message suggesting an infeasible batch. A batch is feasible
under~\eqref{eq:F2} exactly when its operations form a clique in
$\Eset_{OO}^{\mathrm{batch}}$.

\paragraph{Edge features.}
\begin{equation}
  \mathbf{a}_{oo'}=\Bigl[\
  \tfrac{|\kappa(o)-\kappa(o')|}{\Delta},\
  \tfrac{|\Mset_o\cap\Mset_{o'}|}{|\Mset_o\cup\Mset_{o'}|},\
  \tfrac{1}{|\Mset_o\cap\Mset_{o'}|}\textstyle\sum_{\mu\in\Mset_o\cap\Mset_{o'}}\max(p_{o,\mu},p_{o',\mu}),\
  B_{\varphi(o)},\ 0\ \Bigr]\in\R^5.
\end{equation}
The components are the order distance, the machine overlap (Jaccard index), the
processing time of the pair if batched in parallel (the longer of the two), and
the capacity. A zero is appended so that all attributed edge types have the same
dimension. The family has no node of its own in this representation, so its
batch-level parameters are also appended to the operation features:
\begin{equation}
  \mathbf{x}''_o=\Bigl[\ \mathbf{x}'_o \cat B_{\varphi(o)} \cat \bar p_{\varphi(o)}\ \Bigr].
\end{equation}

\paragraph{Size.}
$|\Eset_{OO}^{\mathrm{batch}}|\le\sum_f|\Oset^f|\,(|\Oset^f|-1)$, which is
quadratic in the family size. By \eqref{eq:F1}, $|\Oset^f|\le n$, so the relation
has at most $|\Fset|\,n^2$ edges.

% ===========================================================================
\section{Graph attention message passing}
\label{sec:gat}

Both representations use the same heterogeneous GAT layer. For every relation
$r=(\tau_s,\rho,\tau_d)$ and edge $(u,v)$ of type $r$, the layer computes
(GATv2 with edge attributes):
\begin{align}
  e^{r}_{uv} &= \mathbf{a}_r^{\top}\,\mathrm{LeakyReLU}\bigl(\mathbf{W}^{r}_{s}\mathbf{h}_u+\mathbf{W}^{r}_{d}\mathbf{h}_v+\mathbf{W}^{r}_{e}\mathbf{a}_{uv}\bigr), \label{eq:gat1}\\
  \alpha^{r}_{uv} &= \frac{\exp(e^{r}_{uv})}{\sum_{w\in\mathcal{N}_r(v)}\exp(e^{r}_{wv})}, \label{eq:gat2}\\
  \mathbf{h}^{r}_{v} &= \sum_{u\in\mathcal{N}_r(v)}\alpha^{r}_{uv}\,\mathbf{W}^{r}_{s}\mathbf{h}_u, \label{eq:gat3}\\
  \mathbf{h}^{(l+1)}_{v} &= \sigma\Bigl(\operatorname*{AGG}_{r:\ \tau_d(r)=\operatorname{type}(v)}\mathbf{h}^{r}_{v}\Bigr), \label{eq:gat4}
\end{align}
where the term $\mathbf{W}^{r}_{e}\mathbf{a}_{uv}$ is omitted for relations
without edge attributes, AGG is a sum or mean over relations, and multi-head
outputs are concatenated.

\begin{itemize}
  \item \textbf{Representation A:} information reaches a machine from an operation in two hops (operation $\to$ family $\to$ machine). The attention $\alpha^{\textsf{has}}_{of}$ lets the family node weigh its candidate members, and machines receive batch-level times through $\mathbf{a}_{f\mu}$.
  \item \textbf{Representation B:} batching information is exchanged in one hop (operation $\leftrightarrow$ operation). The attention coefficient $\alpha^{\textsf{batch}}_{oo'}$ directly scores how suitable it is to batch $o$ with $o'$.
\end{itemize}

% ===========================================================================
\section{Consequences for the decision process}
\label{sec:mdp}

Both representations keep the action space of the operation--job--machine graph.
The policy scores the $(\text{machine},\textsf{exec},\text{job})$ edges and selects
one pair $a_t=(\mu,j)$. The two representations therefore differ \emph{only} in
the graph given to the network, which makes their comparison fair and lets the
same agent be used for both.

The selected pair opens a batch $b_t$ for the current operation $o=o_j^t$ on
machine $\mu$, starting at
\begin{equation}
  S_{b_t}=\max\bigl(\text{free time of }\mu,\ \text{ready time of } j\bigr).
\end{equation}
If $o$ belongs to a family $f$, the environment completes the batch with other
jobs $j'$ whose current operation $o'=o_{j'}^t$ satisfies
\begin{equation}
  \varphi(o')=f,\quad \mu\in\Mset_{o'},\quad \text{ready time of } j'\le S_{b_t},
\end{equation}
adding them in order of earliest ready time while
\begin{equation}
  |b_t|\le B_f,\qquad \max_{o\in b_t}\kappa(o)-\min_{o\in b_t}\kappa(o)\le\Delta .
\end{equation}
A candidate joins only if it is already ready, so joining never delays the batch.
The batch then occupies $\mu$ during
$[S_{b_t},\ S_{b_t}+\max_{o\in b_t}p_{o,\mu}]$ (parallel batching), and all of
its jobs advance to their next operation.

\paragraph{Action-edge features.} After every decision each
$(\text{machine},\textsf{exec},\text{job})$ edge $(\mu,j)$ is rebuilt with
\begin{equation}
  \mathbf{a}_{\mu j}=\bigl[\ p_{b},\ \tfrac{p_{o_j^t,\mu}}{\sum_{\mu'}p_{o_j^t,\mu'}},\ s_{j\mu},\ s_{j\mu}+p_{b},\ s_{j\mu}-\text{free time of }\mu\ \bigr],
\end{equation}
where $s_{j\mu}$ is the earliest start of $o_j^t$ on $\mu$. In Representations~A
and~B, $p_b$ is the processing time of the batch the action would dispatch (its
longest member), i.e.\ what the action actually costs. The baseline uses the
individual time $p_{o_j^t,\mu}$, so its graph carries no batching information.

\paragraph{Warm start.} The mask and the dispatch rule used as the warm-start
teacher are those of the operation--job--machine graph: earliest start time, with
ties broken by the earliest completion time. Both quantities are visible to the
network as columns of $\mathbf{a}_{\mu j}$, so the teacher can actually be imitated.

% ===========================================================================
\section{Comparison}
\label{sec:comparison}

To measure what each representation contributes, both are compared with a
\emph{baseline}. The baseline solves the same batching problem with the same
action space and batch completion rule (Section~\ref{sec:mdp}), but on the
unchanged operation--job--machine graph of Section~\ref{sec:base}. It has no
family node, no batch edges and none of the family features of
Section~\ref{sec:familyfeature}, so the network cannot see which operations can
be batched. Any improvement of Representation~A or~B over the baseline can
therefore be attributed to how the batching information is represented.

\begin{table}[h]
\centering
\caption{Comparison of the two representations.}
\label{tab:comparison}
\begin{tabular}{l p{5.2cm} p{5.2cm}}
\toprule
 & A: family node & B: batch edge \\
\midrule
New types        & 1 node type, 2 relations (+ reverses) & 1 relation \\
Number of edges  & $\mathcal{O}\bigl(|\Oset|+\sum_f|\Mset_f|\bigr)$ (linear) & $\mathcal{O}\bigl(\sum_f|\Oset^f|^2\bigr)$ (quadratic in family size) \\
Order difference $\Delta$ & visible through the order spread feature & built into the edge set \\
Batch parameters & stored once on the family node and its edges & duplicated on operation nodes and edges \\
Hops op $\to$ machine & 2 (via family) & 1 (via \textsf{exec}) \\
Action & (machine, job) + batch completion rule & (machine, job) + batch completion rule \\
\bottomrule
\end{tabular}
\end{table}

% ===========================================================================
\section{Illustrative example}
\label{sec:example}

Consider $n=3$ jobs with four operations each, $m=2$ machines, $B_1=B_2=2$ and
$\Delta=2$. Two families are shown:
\[
  \Oset^1=\{O_{1,1},\,O_{2,2},\,O_{3,4}\},\qquad
  \Oset^2=\{O_{1,3},\,O_{2,3}\}.
\]
The remaining operations are singleton families. Assume all operations in these
families are eligible on both machines, and let $p_{O_{1,1},1}=5$ and $p_{O_{2,2},1}=8$.
\begin{itemize}
  \item \eqref{eq:F1} holds: each family has at most one operation per job.
  \item Order difference in $\Oset^1$: $|1-2|=1\le2$ and $|2-4|=2\le2$ are compatible, but $|1-4|=3>2$, so $O_{1,1}$ and $O_{3,4}$ can never share a batch. By~\eqref{eq:B2a}, the feasible batches of size two are $\{O_{1,1},O_{2,2}\}$ and $\{O_{2,2},O_{3,4}\}$.
  \item In $\Oset^2$, $|3-3|=0$, so $\{O_{1,3},O_{2,3}\}$ is feasible.
  \item Parallel batching: if $\{O_{1,1},O_{2,2}\}$ runs on machine 1, then $p_b=\max(5,8)=8$ and both operations finish at $S_b+8$.
  \item Representation A adds two family nodes and $|\Eset_{OF}|=5$ edges, one per operation. It also adds edges to the two machines from each family.
  \item Representation B adds the directed edges $O_{1,1}\leftrightarrow O_{2,2}$ and $O_{2,2}\leftrightarrow O_{3,4}$ in family 1 (not $O_{1,1}\leftrightarrow O_{3,4}$), plus $O_{1,3}\leftrightarrow O_{2,3}$ in family 2. That is 6 edges, excluding self-loops.
\end{itemize}

% ===========================================================================
\section{Diagnostic of the batch process and version 2}
\label{sec:diagnostic}

\subsection{Diagnostic}
The decision process of Section~\ref{sec:mdp} was checked on the 40 instances of the
batching dataset, with the configuration used for training ($\texttt{sel\_k}=1$,
mask on the earliest completion), replaying the warm-start teacher and a random policy
(about 3\,500 decisions each).

\paragraph{Correct behaviour.}
\begin{itemize}
  \item Every batch satisfies the rules: one family, at most $B_f$ members,
    order difference $\le\Delta$, the machine eligible for every member, every member
    ready by the start, one operation per job. All schedules pass the independent checker.
  \item The start time and batch processing time shown on the chosen action edge are
    exactly those the environment then schedules (all 3\,515 decisions, both
    representations).
  \item The partner count $\pi_j(t)$ on the job node equals the number of partners the
    batch rule actually adds in 99.9\% of the decisions.
  \item A partner can lengthen the batch (54\% of the batches, $+15\%$ of the chosen
    job's time on average), since the batch lasts as long as its longest member. This is
    not a defect: only adding partners that do not lengthen it raised the teacher's gap
    to CP-SAT from 0.458 to 0.661 and removed almost all batching. The action edge already
    shows the batch time, so the policy sees the cost.
\end{itemize}

\paragraph{Limitation 1: the batch never waits.}
A partner joins only if it is already ready at the start. The teacher missed 214 compatible
partners that became ready shortly afterwards (median 0.44 batch times later) and batched
11.7\% of the operations, against 18\% in the CP-SAT solutions. Letting the batch wait for
partners ready within half the chosen job's time improved the teacher's mean gap from
0.458 to 0.424, but it helped 18 instances and hurt 17: waiting should be a decision, not
a fixed rule.

\paragraph{Limitation 2: the mask ignores batching.}
The mask ranks each job's machines by the job's \emph{individual} completion. With
$\texttt{sel\_k}=1$, in 14\% of the family decisions a larger batch was possible on
another machine, and that machine was always masked. The action-edge features and the
teacher, however, use the completion of the batch.

\subsection{Version 2}
Version 2 (representations \texttt{ojmb\_node\_v2}, \texttt{ojmb\_edge\_v2},
\texttt{ojmb\_base\_v2}) changes only the action space; the graphs, features and the batch
rule are those of Sections~\ref{sec:familyfeature}--\ref{sec:repB}.

\paragraph{Wait actions.} For an allowed action $(\mu,j)$ whose current operation belongs
to a family, let $b$ be the batch it dispatches now at $s_{j\mu}$, and $s^{w}_{j\mu}$ the
ready time of the next compatible partner: the smallest $r_{j'}>s_{j\mu}$ over the jobs
$j'\notin b$ whose current operation is in the same family, is eligible on $\mu$ and keeps
the order difference of $b\cup\{o_{j'}^t\}$ within $\Delta$ (defined when $|b|<B_f$).
If the batch dispatched at $s^{w}_{j\mu}$ is larger than $b$, a second action edge
$(\mu,j)$ is added, ``wait for the next partner'': it starts at $s^{w}_{j\mu}$ and applies
the same batch rule there. Its features are those of the immediate action with the later
start, the batch time of the larger batch (the individual time for the baseline) and the
machine idle time $s^{w}_{j\mu}-F_\mu$ that waiting leaves. The policy therefore sees both
options and their costs, and chooses.

\paragraph{Batch-aware mask.} Each job keeps the $\texttt{sel\_k}$ machines with the
earliest completion of the batch the action dispatches, $s_{j\mu}+p_b$, instead of
$s_{j\mu}+p_{o_j^t,\mu}$. A wait action is allowed exactly when its immediate twin is.

\paragraph{Effect on the teacher and on a random policy.} A wait option exists in 43\% of
the teacher's decisions. The teacher, which minimises the chosen job's completion, never
takes one, and its gap changes from 0.458 to 0.466 (family node) and from 0.438 to 0.460
(baseline, whose teacher ranks by individual times while the mask now ranks by batch
times). A random policy batches more with the wait actions (20.8\% $\to$ 23.9\% of the
operations). Whether waiting pays off must therefore be learned by BOPO beyond the
teacher; the comparison with version 1 is made on trained policies.

\paragraph{Verification.} The version-1 environments produce exactly the same schedules
before and after the change (12 configurations: three representations, $\texttt{sel\_k}\in
\{1,100\}$, teacher and random policy). \texttt{tests/test\_batching\_v2.py} checks that the
version-2 schedules are feasible, that every wait action is the twin of an allowed action,
starts later, shows its start, and when taken schedules the larger batch at exactly that
start; that the mask keeps the machine with the earliest batch completion; and that the
three representations train.

\end{document}
